% Section 3: Background and Problem Setup
% Note: eq:ii and eq:gradient are the canonical definitions.
% Full formulation details in Appendix (app:formulation-details).

\section{Background and Problem Setup}
\label{sec:background}

We formalize compound sparsity, define the Interaction Index, and introduce phase segmentation. Full mathematical details appear in Appendix~\ref{app:formulation-details}.

\paragraph{Compound sparsity.}
Let $s \in (0,1)$ denote a unified sparsity budget. \emph{Attention head pruning} removes the bottom-$s$ fraction of heads per layer by Frobenius norm importance \cite{voita-acl2019, michel-neurips2019}. \emph{KV cache eviction} removes the bottom-$s$ fraction of token positions from the cache. We compare two strategies: \emph{importance-based}, which scores positions by position-normalized cumulative attention mass (Eq.~\ref{eq:kv-importance} in Appendix~\ref{app:formulation-details}), and \emph{random}, which evicts uniformly at random \cite{nacl-acl2024}. Applying both simultaneously yields \emph{compound sparsity}; the two axes interact because head pruning alters the importance ranking used for eviction \cite{spark-aaai2026}. Existing joint-pruning approaches such as GRASPrune \cite{graspune-acl2026} learn global budget-constrained gates but treat decisions as static across generation.

\paragraph{Interaction Index.}
For a text segment, let $\text{PPL}_c$ denote perplexity under condition $c \in \{d, a, k, b\}$ (dense, attention-only, KV-only, compound). Normalizing by the dense baseline yields $r_c = \text{PPL}_c / \text{PPL}_d$, and:
\begin{equation}
    \mathrm{II} = r_b - r_a - r_k + 1.
    \label{eq:ii}
\end{equation}
$\mathrm{II} > 0$ indicates super-additive degradation; $\mathrm{II} = 0$ indicates additive independence. We adopt the additive null as primary: it assumes $r_b = r_a + r_k - 1$ (compound loss equals sum of individual losses), the most conservative independence assumption. Although the two axes are mechanistically linked (\S\ref{sec:analysis}), rejection ($\mathrm{II} > 0$) establishes super-additivity regardless of null specification. Appendix~\ref{app:multiplicative-null} reports parallel results under the multiplicative null ($r_b = r_a \cdot r_k$); the phase-severity gradient is preserved ($G_{\times}{=}21$--$64{\times}$), supporting null-choice invariance. The II distribution is heavily right-skewed; we report medians throughout.

\paragraph{Severity gradient.}
To capture phase variation, the per-problem \emph{amplification ratio} is:
\begin{equation}
    G^{(i)} = \frac{\mathrm{II}^{(i)}_{\,\textit{low}}}{\mathrm{II}^{(i)}_{\,\textit{high}}},
    \label{eq:gradient}
\end{equation}
where $\mathrm{II}^{(i)}_{\,\textit{low/high}}$ is the interaction index in the low/high-sparsity phase for problem $i$. We report $G = \widetilde{G^{(i)}}$, the median across problems, with bootstrap CIs. We use the median of per-problem ratios (not the ratio of medians) to preserve per-problem pairing; see Appendix~\ref{app:profiling-subset} for comparison. $G$ serves as an auxiliary descriptive ratio that illustrates the degree of phase concentration; the primary statistical evidence for the severity gradient rests on $\Delta$ and paired Wilcoxon tests (\S\ref{sec:phase-amplification}).

\paragraph{Phase segmentation.}
\label{par:phase-segmentation}
Inspired by ThinKV's phase-aware perspective on reasoning-time compression~\cite{thinkv-iclr2026}, each generated sequence is partitioned into three phases (low-, mid-, high-sparsity) by tercile boundaries of per-position attention sparsity, using a quantile-based segmentation rather than explicit reasoning delimiters (e.g., \texttt{<think>} tokens). Low-sparsity phases correspond to broadly distributed attention (exploratory reasoning); high-sparsity phases to concentrated attention (structured computation). We note that this segmentation is an \emph{operationalization choice}: attention sparsity serves as a continuous proxy for reasoning phase rather than a theoretically grounded boundary. We validate robustness to segmentation granularity (tercile/quartile/quintile; Appendix~\ref{app:seg-robust}) and minimum phase length (Appendix~\ref{app:mpl-sensitivity}). Validation against DeepSeek's explicit \texttt{<think>} markers using adaptive per-sequence quantile segmentation yields substantial agreement ($\kappa{=}0.88$, AUC-ROC${=}0.98$; $n{=}48$; Appendix~\ref{app:think-validation}). Importantly, phases reflect genuine computational mode differences even under the dense (uncompressed) baseline: per-phase dense perplexity increases monotonically from high ($1.10$) through mid ($1.17$) to low ($1.18$), confirming that low-sparsity phases correspond to inherently harder computation; the phase gradient in compound severity (\S\ref{sec:results}) thus amplifies a pre-existing computational asymmetry rather than creating one. Details in Appendix~\ref{app:formulation-details}.
